\documentclass[12pt,a4paper]{article}

\usepackage[utf8]{inputenc}
\usepackage[T1]{fontenc}
\usepackage{lmodern}
\usepackage{geometry}
\usepackage{setspace}
\usepackage{hyperref}
\usepackage{longtable}
\usepackage{booktabs}
\usepackage{array}
\usepackage{xcolor}
\usepackage{listings}
\usepackage{caption}
\usepackage{float}
\usepackage{enumitem}

\geometry{
    left=25mm,
    right=25mm,
    top=25mm,
    bottom=25mm
}

\hypersetup{
    colorlinks=true,
    linkcolor=blue,
    urlcolor=blue,
    citecolor=blue,
    pdftitle={Network Intrusion Detection System Using Machine Learning},
    pdfauthor={NIDS Project}
}

\onehalfspacing
\setlength{\parindent}{0pt}
\setlength{\parskip}{6pt}

\lstdefinestyle{codeblock}{
    basicstyle=\ttfamily\small,
    backgroundcolor=\color{gray!10},
    frame=single,
    breaklines=true,
    columns=fullflexible
}

\lstset{style=codeblock}

\begin{document}

\begin{titlepage}
    \centering
    \vspace*{2cm}
    {\Huge \textbf{Network Intrusion Detection System}}\\[0.4cm]
    {\LARGE \textbf{Using Machine Learning}}\\[1.5cm]
    {\Large Detailed Project Report}\\[1.5cm]
    \vfill
    {\large Dataset: CICIDS2017}\\[0.3cm]
    {\large Models: Logistic Regression, KNN, Random Forest, XGBoost}\\[0.3cm]
    {\large Inference: Offline PCAP Flow Analysis}\\[2cm]
    \vfill
    {\large \today}
\end{titlepage}

\pagenumbering{roman}

\section*{Abstract}
\addcontentsline{toc}{section}{Abstract}

This project develops a Machine Learning-based Network Intrusion Detection System (NIDS) for classifying network traffic flows into normal and malicious categories. The system is built around the CICIDS2017 dataset, a widely used benchmark dataset for intrusion detection research. The project implements a full training and evaluation pipeline, including dataset loading, label grouping, data cleaning, feature scaling, class imbalance handling, model training, model comparison, model persistence, and offline PCAP-based inference.

The main task is multi-class classification. The final label space includes \texttt{Normal}, \texttt{DoS}, \texttt{DDoS}, \texttt{PortScan}, \texttt{BruteForce}, \texttt{Bot}, and \texttt{Infiltration}. Four machine learning models are evaluated: Logistic Regression, K-Nearest Neighbors, Random Forest, and XGBoost. Because CICIDS2017 is highly imbalanced, especially for minority classes such as \texttt{Bot} and \texttt{Infiltration}, the project also evaluates model behavior before and after conservative imbalance handling. The imbalance strategy combines Random Under-Sampling for the majority \texttt{Normal} class with light SMOTE for selected minority classes.

Experimental results show that Random Forest without imbalance handling achieves the best full-dataset Test Macro-F1 score of 0.9378, while XGBoost with conservative imbalance handling shows the clearest improvement on minority-class detection, increasing Test Macro-F1 from 0.8243 to 0.9215 and improving \texttt{Infiltration} F1 from 0.00 to 0.71. Logistic Regression performs poorly on rare attack classes, while KNN achieves moderate full-dataset performance but is computationally expensive during inference.

The project also supports offline PCAP testing. PCAP files are converted into CICIDS-like flow features using a patched CICFlowMeter-based extractor or a Scapy fallback. The trained model then predicts traffic classes from the extracted flows. This inference workflow demonstrates how the trained NIDS can be applied beyond static CSV evaluation, while also showing the practical limitations caused by domain shift, feature extraction differences, and the mismatch between benchmark flow data and external packet captures.

\clearpage
\tableofcontents
\clearpage
\pagenumbering{arabic}

\section{Introduction}

\subsection{Background}

Network Intrusion Detection Systems monitor network traffic and detect suspicious or malicious activity. Traditional intrusion detection systems often rely on signature-based rules. Although signature-based methods can detect known attacks reliably, they usually struggle with unseen attack patterns, modified payloads, and zero-day behavior.

Machine learning offers a different approach. Instead of checking only fixed rules, a model learns statistical patterns from historical network traffic. Once trained, it can classify new traffic based on flow-level behavior such as packet counts, byte counts, packet lengths, inter-arrival times, protocol flags, and traffic direction statistics.

This project focuses on flow-based intrusion detection. A flow represents a communication session or a group of packets sharing common properties such as source, destination, ports, and protocol. Flow-based features are more compact than raw packets and are commonly used in intrusion detection datasets such as CICIDS2017.

\subsection{Project Motivation}

The goal of this project is not only to train a high-accuracy model, but also to understand model behavior under dataset imbalance and evaluate whether trained models can be used on offline PCAP traffic.

The project investigates the following questions:

\begin{itemize}
    \item Can machine learning classify common network attacks from CICIDS2017 flow features?
    \item How does class imbalance affect the result?
    \item Does imbalance handling improve all models equally?
    \item Which model is most suitable for benchmark reporting?
    \item Which model is most useful for PCAP inference?
    \item Can a model trained on CICIDS2017 generalize to external PCAP files?
\end{itemize}

\subsection{Project Scope}

The project includes CICIDS2017 CSV dataset loading, multi-class label grouping, data cleaning, preprocessing, train/validation/test splitting, conservative dataset imbalance handling, model training, model comparison, artifact saving, and offline PCAP inference through flow extraction.

The project no longer includes rule-based log analysis. The previous log analysis workflow was removed to keep the project focused on machine learning and PCAP-based NIDS inference.

\section{Project Structure}

The project is organized into a small number of main directories:

\begin{lstlisting}
NIDS---Network-Intrusion-Detection-System/
|-- data/
|   |-- raw/CICIDS2017/
|   `-- pcap/
|-- models/final/
|-- experiments/models/
|-- outputs/
|-- reports/training_runs/
|-- src/
|   |-- data/
|   |-- models/
|   |-- pipeline/
|   `-- utils/
|-- live_analyze.py
|-- main.py
|-- requirements.txt
|-- README.md
`-- REPORT.md
\end{lstlisting}

\subsection{Data Directory}

The \texttt{data/raw/CICIDS2017/} folder contains the CICIDS2017 CSV flow files used for training and evaluation. The \texttt{data/pcap/} folder contains offline PCAP files used for inference testing.

\subsection{Model Directories}

The \texttt{models/final/} directory is the default inference directory. It is generated locally after training and is intentionally not committed to Git because model artifacts can be large and environment-specific. It contains the model artifacts used by PCAP inference:

\begin{lstlisting}
model.pkl
scaler.pkl
label_encoder.pkl
features.pkl
\end{lstlisting}

The \texttt{experiments/models/} folder stores artifacts from different experiment runs. It is separated from \texttt{models/final/} to keep the main model directory clean. These experiment artifacts are also local runtime outputs and are not required in the GitHub repository.

\subsection{Source Code}

The \texttt{src/} directory contains reusable source code. The \texttt{src/data/} package handles loading, cleaning, label mapping, and imbalance handling. The \texttt{src/models/} package contains model definitions and prediction helpers. The \texttt{src/pipeline/} package contains the end-to-end training pipeline.

\section{Dataset}

\subsection{Dataset Source}

The project uses the CICIDS2017 flow dataset. The dataset is stored in:

\begin{lstlisting}
data/raw/CICIDS2017/
\end{lstlisting}

The loaded dataset contains 2,830,743 rows and 79 columns including the label column.

\subsection{CSV Files Used}

The full dataset is loaded from eight CICIDS2017 CSV files:

\begin{itemize}
    \item \texttt{Monday-WorkingHours.pcap\_ISCX.csv}
    \item \texttt{Tuesday-WorkingHours.pcap\_ISCX.csv}
    \item \texttt{Wednesday-workingHours.pcap\_ISCX.csv}
    \item \texttt{Thursday-WorkingHours-Morning-WebAttacks.pcap\_ISCX.csv}
    \item \texttt{Thursday-WorkingHours-Afternoon-Infilteration.pcap\_ISCX.csv}
    \item \texttt{Friday-WorkingHours-Morning.pcap\_ISCX.csv}
    \item \texttt{Friday-WorkingHours-Afternoon-PortScan.pcap\_ISCX.csv}
    \item \texttt{Friday-WorkingHours-Afternoon-DDos.pcap\_ISCX.csv}
\end{itemize}

\subsection{Feature Type}

The dataset uses flow-level features rather than raw packets. Example feature groups include flow duration, forward and backward packet counts, byte statistics, packet length statistics, flow rate, inter-arrival time statistics, TCP flag counts, header length statistics, and active/idle time statistics.

\subsection{Original Class Distribution}

\begin{longtable}{l r}
\caption{Original class distribution after label grouping.}\\
\toprule
\textbf{Class} & \textbf{Samples} \\
\midrule
\endfirsthead
\toprule
\textbf{Class} & \textbf{Samples} \\
\midrule
\endhead
Normal & 2,273,097 \\
DoS & 252,661 \\
PortScan & 158,930 \\
DDoS & 128,027 \\
BruteForce & 13,835 \\
Bot & 1,966 \\
Infiltration & 47 \\
\bottomrule
\end{longtable}

\subsection{Dataset Imbalance}

The dataset is highly imbalanced. \texttt{Normal} traffic dominates the dataset, while \texttt{Infiltration} has only 47 samples. Accuracy and Weighted-F1 can therefore be misleading. A model can perform very well on majority classes while failing rare attacks. For this reason, Macro-F1 is emphasized throughout the evaluation.

\section{Label Grouping}

The raw CICIDS2017 labels are mapped into broader attack categories.

\begin{longtable}{p{0.48\textwidth} p{0.35\textwidth}}
\caption{Label grouping used in preprocessing.}\\
\toprule
\textbf{Original Label} & \textbf{Grouped Label} \\
\midrule
\endfirsthead
\toprule
\textbf{Original Label} & \textbf{Grouped Label} \\
\midrule
\endhead
\texttt{BENIGN}, \texttt{Benign}, \texttt{Normal} & \texttt{Normal} \\
\texttt{DoS Hulk}, \texttt{DoS GoldenEye}, \texttt{DoS Slowloris}, \texttt{DoS Slowhttptest}, \texttt{DoS} & \texttt{DoS} \\
\texttt{DDoS} & \texttt{DDoS} \\
\texttt{PortScan} & \texttt{PortScan} \\
\texttt{FTP-Patator}, \texttt{SSH-Patator}, \texttt{BruteForce} & \texttt{BruteForce} \\
\texttt{Bot} & \texttt{Bot} \\
\texttt{Web Attack Brute Force}, \texttt{Web Attack XSS}, \texttt{Web Attack Sql Injection} & \texttt{Web} \\
\texttt{Infiltration}, \texttt{Heartbleed} & \texttt{Infiltration} \\
\bottomrule
\end{longtable}

In the final full-dataset experiments, the observed classes are \texttt{Normal}, \texttt{DoS}, \texttt{DDoS}, \texttt{PortScan}, \texttt{BruteForce}, \texttt{Bot}, and \texttt{Infiltration}.

\section{Preprocessing Pipeline}

\subsection{Column Name Normalization}

Column names are stripped to remove leading or trailing spaces:

\begin{lstlisting}[language=Python]
df_clean.columns = df_clean.columns.str.strip()
\end{lstlisting}

\subsection{Metadata Removal}

Metadata columns are removed before training:

\begin{lstlisting}
Flow ID
Source IP
Source Port
Destination IP
Destination Port
Protocol
Timestamp
\end{lstlisting}

Equivalent fields from PCAP flow extractors are also removed, such as \texttt{flow\_id}, \texttt{src\_ip}, \texttt{src\_port}, \texttt{dst\_ip}, \texttt{dst\_port}, \texttt{protocol}, and \texttt{timestamp}. These fields identify traffic endpoints rather than general behavior and could cause the model to learn dataset-specific identifiers.

\subsection{Missing and Infinite Values}

Invalid numerical values are handled as follows:

\begin{lstlisting}[language=Python]
df_clean.replace([np.inf, -np.inf], np.nan, inplace=True)
df_clean.fillna(0, inplace=True)
\end{lstlisting}

\subsection{Numeric Conversion}

All feature columns are converted to numeric values. Failed conversions become \texttt{NaN} and are later filled with zero.

\subsection{Train/Validation/Test Split}

The cleaned dataset is split into train, validation, and test subsets.

\begin{longtable}{l r}
\caption{Dataset split ratio.}\\
\toprule
\textbf{Split} & \textbf{Ratio} \\
\midrule
Train & 70\% \\
Validation & 10\% \\
Test & 20\% \\
\bottomrule
\end{longtable}

The split is stratified when possible so that class proportions are preserved.

\subsection{Label Encoding and Feature Scaling}

Text labels are encoded using \texttt{LabelEncoder}. Features are scaled using \texttt{StandardScaler}. The scaler is fitted only on the training set, while validation, test, and inference data only use \texttt{transform}. This avoids data leakage.

\section{Class Imbalance Handling}

\subsection{Strategy}

The imbalance handling function uses a conservative strategy:

\begin{enumerate}
    \item Random under-sampling for the majority \texttt{Normal} class.
    \item Light SMOTE for selected minority classes.
    \item No aggressive oversampling for extremely rare classes.
\end{enumerate}

\subsection{Normal Under-Sampling}

The \texttt{Normal} class is reduced to at most 300,000 training samples:

\begin{lstlisting}[language=Python]
target_normal = min(300000, normal_count)
\end{lstlisting}

This reduces majority-class dominance while still keeping a large number of real normal samples.

\subsection{Conservative SMOTE}

\begin{longtable}{p{0.45\textwidth} p{0.45\textwidth}}
\caption{Conservative SMOTE strategy.}\\
\toprule
\textbf{Condition} & \textbf{Strategy} \\
\midrule
\endfirsthead
\toprule
\textbf{Condition} & \textbf{Strategy} \\
\midrule
\endhead
Class has fewer than 100 samples & Do not increase \\
\texttt{Bot} & Increase up to 1.5x, capped at 2,500 \\
Class has fewer than 2,000 samples & Increase up to 1.5x, capped at 4,000 \\
Class has fewer than 10,000 samples & Increase up to 1.25x, capped at 12,000 \\
Large classes & Keep unchanged \\
\bottomrule
\end{longtable}

\subsection{Training Distribution After Imbalance Handling}

\begin{longtable}{l r}
\caption{Training distribution after imbalance handling.}\\
\toprule
\textbf{Class} & \textbf{Training Samples} \\
\midrule
Normal & 300,000 \\
DoS & 176,863 \\
PortScan & 111,251 \\
DDoS & 89,618 \\
BruteForce & 12,000 \\
Bot & 2,064 \\
Infiltration & 33 \\
\bottomrule
\end{longtable}

The original \texttt{Infiltration} count is 47 in the full dataset. The value 33 appears after the train/validation/test split because imbalance handling is applied only to the training set.

\section{Machine Learning Models}

\subsection{Logistic Regression}

Logistic Regression is used as a linear baseline.

\begin{lstlisting}
solver = saga
max_iter = 300
tol = 1e-3
n_jobs = -1
random_state = 42
\end{lstlisting}

It performs poorly on rare attack classes and produces convergence warnings on the full dataset.

\subsection{K-Nearest Neighbors}

KNN is used as a distance-based baseline.

\begin{lstlisting}
n_neighbors = 5
n_jobs = -1
\end{lstlisting}

KNN achieves moderate Macro-F1 on the full dataset, but inference is slow because prediction requires distance comparisons against a large training set.

\subsection{Random Forest}

Random Forest is a tree-based ensemble model.

\begin{lstlisting}
n_estimators = 100
n_jobs = -1
random_state = 42
\end{lstlisting}

It achieves the best full-dataset Test Macro-F1 on the internal CICIDS2017 split.

\subsection{XGBoost}

XGBoost is the main boosted tree model.

\begin{lstlisting}
n_estimators = 600
max_depth = 4
learning_rate = 0.04
subsample = 0.8
colsample_bytree = 0.75
reg_alpha = 2.0
reg_lambda = 8.0
gamma = 2.0
min_child_weight = 10
max_delta_step = 1
objective = multi:softmax
eval_metric = merror
tree_method = hist
n_jobs = -1
random_state = 42
\end{lstlisting}

When XGBoost is trained with imbalance handling, clipped balanced sample weights are also used:

\begin{lstlisting}
min weight = 0.75
max weight = 2.00
\end{lstlisting}

XGBoost shows the clearest improvement after imbalance handling and performs best in PCAP model-only testing.

\section{Training Pipeline}

The full training process is:

\begin{lstlisting}
Load CICIDS2017 CSV files
Clean data and map labels
Split into train/validation/test
Apply imbalance handling if enabled
Encode labels
Scale features
Train selected model
Evaluate validation set
Evaluate test set
Run sampled cross-validation if possible
Save model artifacts
\end{lstlisting}

\subsection{Training Commands}

Train XGBoost with imbalance handling:

\begin{lstlisting}
python3 main.py --train data/raw/CICIDS2017 --model xgb
\end{lstlisting}

Train Random Forest:

\begin{lstlisting}
python3 main.py --train data/raw/CICIDS2017 --model rf
\end{lstlisting}

Train without imbalance handling:

\begin{lstlisting}
python3 main.py --train data/raw/CICIDS2017 --model rf --no-imbalance
\end{lstlisting}

Train and save artifacts to a custom directory:

\begin{lstlisting}
python3 main.py --train data/raw/CICIDS2017 --model xgb --model-dir models/final
\end{lstlisting}

\section{Evaluation Metrics}

The project uses Precision, Recall, F1-score, Macro-F1, Weighted-F1, and sampled cross-validation where applicable.

\subsection{Macro-F1}

Macro-F1 averages F1-score across all classes equally. It is the most important metric for this project because it reveals whether minority classes are being detected.

\subsection{Weighted-F1}

Weighted-F1 averages F1-score by class support. Because \texttt{Normal} has very high support, Weighted-F1 can remain high even if rare classes perform poorly.

\section{Full Dataset Results}

\subsection{Overall Results}

\begin{longtable}{l l r r r r l}
\caption{Full-dataset validation and test results.}\\
\toprule
\textbf{Model} & \textbf{Imbalance} & \textbf{Val Macro-F1} & \textbf{Val Weighted-F1} & \textbf{Test Macro-F1} & \textbf{Test Weighted-F1} & \textbf{CV Weighted-F1} \\
\midrule
\endfirsthead
\toprule
\textbf{Model} & \textbf{Imbalance} & \textbf{Val Macro-F1} & \textbf{Val Weighted-F1} & \textbf{Test Macro-F1} & \textbf{Test Weighted-F1} & \textbf{CV Weighted-F1} \\
\midrule
\endhead
Logistic Regression & No & 0.4788 & 0.9224 & 0.4791 & 0.9224 & Skipped \\
Logistic Regression & Yes & 0.5030 & 0.8940 & 0.4554 & 0.8938 & 0.8821 $\pm$ 0.0018 \\
KNN & No & 0.8818 & 0.9981 & 0.8769 & 0.9982 & 0.9834 $\pm$ 0.0011 \\
KNN & Yes & 0.8252 & 0.9916 & 0.8378 & 0.9914 & Skipped \\
XGBoost & No & 0.8682 & 0.9987 & 0.8243 & 0.9988 & Skipped \\
XGBoost & Yes & 0.9459 & 0.9984 & 0.9215 & 0.9985 & 0.9981 $\pm$ 0.0002 \\
Random Forest & No & 0.9536 & 0.9988 & 0.9378 & 0.9987 & Skipped \\
Random Forest & Yes & 0.9497 & 0.9983 & 0.9280 & 0.9985 & 0.9979 $\pm$ 0.0002 \\
\bottomrule
\end{longtable}

\subsection{Per-Class Test F1}

\begin{longtable}{l l r r r r r r r}
\caption{Per-class Test F1 scores.}\\
\toprule
\textbf{Model} & \textbf{Imb.} & \textbf{Bot} & \textbf{BruteForce} & \textbf{DDoS} & \textbf{DoS} & \textbf{Infiltration} & \textbf{Normal} & \textbf{PortScan} \\
\midrule
\endfirsthead
\toprule
\textbf{Model} & \textbf{Imb.} & \textbf{Bot} & \textbf{BruteForce} & \textbf{DDoS} & \textbf{DoS} & \textbf{Infiltration} & \textbf{Normal} & \textbf{PortScan} \\
\midrule
\endhead
Logistic Regression & No & 0.00 & 0.00 & 0.75 & 0.78 & 0.00 & 0.96 & 0.87 \\
Logistic Regression & Yes & 0.00 & 0.00 & 0.69 & 0.76 & 0.00 & 0.93 & 0.80 \\
KNN & No & 0.70 & 0.99 & 1.00 & 1.00 & 0.46 & 1.00 & 1.00 \\
KNN & Yes & 0.55 & 0.91 & 1.00 & 0.99 & 0.46 & 0.99 & 0.96 \\
XGBoost & No & 0.78 & 1.00 & 1.00 & 1.00 & 0.00 & 1.00 & 1.00 \\
XGBoost & Yes & 0.75 & 1.00 & 1.00 & 1.00 & 0.71 & 1.00 & 1.00 \\
Random Forest & No & 0.77 & 1.00 & 1.00 & 1.00 & 0.80 & 1.00 & 1.00 \\
Random Forest & Yes & 0.79 & 1.00 & 1.00 & 1.00 & 0.71 & 1.00 & 1.00 \\
\bottomrule
\end{longtable}

\subsection{XGBoost Confusion Matrix and Diagnostics}

The project also saves additional XGBoost diagnostics. The confusion matrix below corresponds to the XGBoost model with conservative imbalance handling. Rows represent true labels and columns represent predicted labels.

{\scriptsize
\begin{longtable}{l r r r r r r r}
\caption{XGBoost confusion matrix with conservative imbalance handling.}\\
\toprule
\textbf{True / Pred.} & \textbf{Bot} & \textbf{BruteForce} & \textbf{DDoS} & \textbf{DoS} & \textbf{Infiltration} & \textbf{Normal} & \textbf{PortScan} \\
\midrule
\endfirsthead
\toprule
\textbf{True / Pred.} & \textbf{Bot} & \textbf{BruteForce} & \textbf{DDoS} & \textbf{DoS} & \textbf{Infiltration} & \textbf{Normal} & \textbf{PortScan} \\
\midrule
\endhead
Bot & 383 & 0 & 0 & 0 & 0 & 10 & 0 \\
BruteForce & 0 & 2758 & 0 & 1 & 0 & 8 & 0 \\
DDoS & 0 & 0 & 25600 & 0 & 0 & 6 & 0 \\
DoS & 0 & 0 & 0 & 50514 & 0 & 18 & 0 \\
Infiltration & 0 & 0 & 0 & 1 & 5 & 3 & 0 \\
Normal & 247 & 13 & 40 & 308 & 0 & 453804 & 208 \\
PortScan & 0 & 0 & 0 & 7 & 0 & 4 & 31775 \\
\bottomrule
\end{longtable}
}

The confusion matrix shows that most high-support attack classes are classified correctly. The rare \texttt{Infiltration} class remains difficult because the test set contains only 9 samples.

\begin{longtable}{r l r}
\caption{Top XGBoost feature importance values.}\\
\toprule
\textbf{Rank} & \textbf{Feature} & \textbf{Importance} \\
\midrule
1 & Subflow Fwd Packets & 0.0936 \\
2 & Bwd Packets/s & 0.0791 \\
3 & act\_data\_pkt\_fwd & 0.0777 \\
4 & Total Fwd Packets & 0.0688 \\
5 & Bwd Packet Length Std & 0.0609 \\
6 & Min Packet Length & 0.0583 \\
7 & Fwd Packet Length Mean & 0.0435 \\
8 & Avg Fwd Segment Size & 0.0434 \\
9 & Idle Max & 0.0359 \\
10 & Subflow Bwd Bytes & 0.0348 \\
\bottomrule
\end{longtable}

\subsection{Interpretation}

Random Forest without imbalance handling achieves the best internal CICIDS2017 Test Macro-F1 of 0.9378. XGBoost benefits the most from imbalance handling, improving Test Macro-F1 from 0.8243 to 0.9215 and improving \texttt{Infiltration} F1 from 0.00 to 0.71. KNN performs worse after imbalance handling because it is sensitive to sample density. Logistic Regression is the weakest model and should be treated only as a baseline.

\section{Model Selection}

For internal CICIDS2017 evaluation, the strongest model is Random Forest without imbalance handling. For imbalance-focused analysis and PCAP-oriented testing, the preferred model is XGBoost with conservative imbalance handling.

\section{PCAP Inference}

\subsection{Inference Flow}

\begin{lstlisting}
Input PCAP
Extract flow features
Normalize feature columns
Clean data
Align columns with training feature list
Scale using saved scaler
Predict using saved model
Decode labels
Print NIDS report
\end{lstlisting}

\subsection{Default Model Directory}

By default, PCAP testing uses:

\begin{lstlisting}
models/final/
\end{lstlisting}

Required files:

\begin{lstlisting}
model.pkl
scaler.pkl
label_encoder.pkl
features.pkl
\end{lstlisting}

\subsection{PCAP Commands}

Default model:

\begin{lstlisting}
python3 main.py --pcap data/pcap/test1.pcap
\end{lstlisting}

Specific model:

\begin{lstlisting}
python3 main.py --pcap data/pcap/task3.dos_victim.pcap --model-dir experiments/models/xgb_with_imbalance
\end{lstlisting}

\subsection{PCAP Test Results}

\subsubsection{\texttt{task3.dos\_victim.pcap}}

This file extracted 41,830 flows using \texttt{cicflowmeter-patched}.

\begin{longtable}{l r r l}
\caption{PCAP model comparison on task3.dos\_victim.pcap.}\\
\toprule
\textbf{Model} & \textbf{Normal} & \textbf{DoS} & \textbf{Interpretation} \\
\midrule
Logistic Regression & 41,830 & 0 & All flows predicted Normal \\
KNN & 41,830 & 0 & All flows predicted Normal \\
Random Forest & 41,830 & 0 & All flows predicted Normal \\
XGBoost with imbalance & 36,368 & 5,462 & Detected DoS-like flows \\
\bottomrule
\end{longtable}

\subsubsection{Other PCAP Files}

\begin{longtable}{l l l}
\caption{Other PCAP test results.}\\
\toprule
\textbf{PCAP} & \textbf{Extracted Flows} & \textbf{Result} \\
\midrule
\texttt{task1.dos\_victim.pcap} & 224 & All models predicted Normal \\
\texttt{task3.dos\_attacker.pcap} & 43,525 & All models predicted Normal \\
\texttt{task1.dos\_attacker.pcap} & 5,114 & All models predicted Normal \\
\bottomrule
\end{longtable}

\subsection{PCAP Interpretation}

The PCAP results show that high performance on the CICIDS2017 internal split does not guarantee strong generalization to external PCAP traffic. Domain shift and flow extraction differences matter. XGBoost with imbalance handling is the only tested model that detects DoS on \texttt{task3.dos\_victim.pcap}.

\section{Usage Guide}

\subsection{Install Dependencies}

\begin{lstlisting}
pip install -r requirements.txt
\end{lstlisting}

\subsection{Train a Model}

\begin{lstlisting}
python3 main.py --train data/raw/CICIDS2017 --model xgb --model-dir models/final
\end{lstlisting}

\subsection{Test a PCAP File}

\begin{lstlisting}
python3 main.py --pcap data/pcap/test1.pcap
\end{lstlisting}

\subsection{Test With a Specific Model}

\begin{lstlisting}
python3 main.py --pcap data/pcap/task3.dos_victim.pcap --model-dir experiments/models/xgb_with_imbalance
\end{lstlisting}

\section{Limitations}

\begin{itemize}
    \item CICIDS2017 has extremely rare classes, especially \texttt{Infiltration}.
    \item Rare-class metrics can fluctuate heavily due to small test support.
    \item PCAP testing is affected by domain shift.
    \item Flow extraction may not perfectly match the CICIDS2017 feature generation process.
    \item KNN has high inference cost on large datasets.
    \item Logistic Regression does not converge well on the full dataset.
\end{itemize}

\section{Future Work}

Future improvements include adding a dedicated CSV-flow prediction command, improving PCAP feature extraction, adding probability thresholding for XGBoost, adding SHAP-based explainability, separating binary detection from multi-class classification, trying LightGBM or CatBoost, evaluating more external PCAP datasets, and adding clearer experiment tracking.

\section{Conclusion}

This project successfully implements a full Machine Learning-based NIDS pipeline using CICIDS2017. It supports data loading, cleaning, label grouping, imbalance handling, model training, model evaluation, artifact saving, and offline PCAP inference.

The experiments show that Random Forest achieves the best internal CICIDS2017 full-dataset Macro-F1 score. However, XGBoost with conservative imbalance handling provides the strongest evidence that imbalance handling improves rare attack detection, especially for \texttt{Infiltration}. In PCAP inference, XGBoost is also the only tested model that detects DoS-like traffic in \texttt{task3.dos\_victim.pcap}.

The project demonstrates an important real-world lesson: high test performance on a benchmark dataset does not automatically guarantee strong performance on external PCAP traffic. Domain shift and flow extraction differences matter. Therefore, the project should present Random Forest as the strongest internal benchmark model and XGBoost with imbalance handling as the most practical model for minority-class and PCAP-oriented analysis.

\end{document}
